%Erzeugt mit dem LaTeX-Generator: http://latex.sehnot.de

%Schriftgröße, Layout, Papierformat, Art des Dokumentes
\documentclass[12pt,oneside,a4paper]{scrbook}

%Einstellungen der Seitenränder
\usepackage[left=3cm,right=2cm,top=2cm,bottom=2cm,includeheadfoot]{geometry}

%neue Rechtschreibung
%\usepackage[ngerman]{babel}

%Umlaute ermöglichen
% \usepackage{fontspec}
% \setromanfont{TeX Gyre Pagella}

% \usepackage[math-style=iso]{unicode-math}
% \setmathfont{Asana Math}

\usepackage{graphicx}
\usepackage{wrapfig}
\usepackage[utf8x]{inputenc}
\usepackage{url}
\usepackage{multirow}
\usepackage{amsmath}
\usepackage{hyperref}
\usepackage{multicol}
\usepackage{amsthm}
\usepackage[official]{eurosym}
\usepackage{setspace}
\newtheorem{definition}{Definition}
\newtheorem{satz}{Theorem}
\theoremstyle{definition}
%Zitieren
\usepackage[longnamesfirst, authoryear]{natbib}
\usepackage{minted}
\usepackage{booktabs}



%Kopf- und Fußzeile
\usepackage{fancyhdr}
\pagestyle{fancy}
\fancyhf{}

\newcommand{\HRule}{\rule{\linewidth}{0.5mm}}

%Kopfzeile rechts bzw. außen
\fancyhead[R]{\nouppercase{\leftmark}}
%Linie oben
\renewcommand{\headrulewidth}{0.5pt}

%Fußzeile links bzw. innen
\fancyfoot[L]{Units of Measure - a Scala Macro System}
%Fußzeile rechts bzw. außen
\fancyfoot[R]{\thepage}
%Linie unten
\renewcommand{\footrulewidth}{0.5pt}

\clubpenalty = 50000
\widowpenalty = 50000




\title{Units of Measure - a Scala Macro System}
\author{Julian Schrittwieser}
\date{\today}

\begin{document}


\begin{titlepage}

\begin{center}


% Upper part of the page
\includegraphics[width=0.4\textwidth]{tu-logo.png}\\[3cm]


\textsc{\Large Software \& Information Engineering}\\[0.5cm]


% Title
\HRule \\[0.8cm]
{ \huge \bfseries Units of Measure}\\[0.4cm]
\Large A Scala Macro System

\HRule \\[1.5cm]

\vfill


% Author and supervisor
\begin{minipage}{0.4\textwidth}
\begin{flushleft} \large
\emph{Author:}\\
Julian \textsc{Schrittwieser}
\end{flushleft}
\end{minipage}
\begin{minipage}{0.4\textwidth}
\begin{flushright} \large
\emph{Advisor:} \\
Univ.-Prof. Dr. Jens \textsc{Knoop}
\end{flushright}
\end{minipage}

\vfill

% Bottom of the page
{\large Februar 12th, 2013}

\end{center}

\end{titlepage}

\pagenumbering{roman} % Roman numerals, other styles possible: http://www.image.ufl.edu/help/latex/intext.shtml
\setcounter{page}{2}

% \onehalfspacing
\chapter*{Preface}

Lorem ipsum dolor sit amet, consectetur adipiscing elit. Donec in ligula et massa ultrices luctus. Aenean nec orci lectus. Maecenas ut ligula non eros dictum porttitor. Nulla vel libero vitae dolor vehicula luctus. Maecenas venenatis sapien eu mi vehicula non dapibus lorem cursus. In dapibus ornare turpis, id tincidunt odio consectetur sit amet. Integer placerat varius velit at semper. Morbi id erat felis. Donec scelerisque eros in diam sollicitudin dictum. In lobortis imperdiet sapien. Fusce volutpat orci sit amet orci consectetur faucibus.

Nullam mattis felis ac est feugiat porta. Suspendisse potenti. Suspendisse non augue eu dolor ultrices dignissim id ut augue. Donec sit amet mi a leo egestas lobortis sed a lorem. Nam sit amet turpis non mi gravida tempor id iaculis purus. Duis hendrerit justo quis orci pellentesque eu luctus magna facilisis. Nulla convallis dui vitae lorem ultrices dignissim. Praesent interdum volutpat aliquam. Nullam massa felis, accumsan vel congue hendrerit, tristique at lectus. Integer ac ipsum nec nisl elementum consequat ut vel erat. Suspendisse potenti. Proin varius, orci vel imperdiet rhoncus, nisi nibh lobortis erat, eget ultrices nibh tortor ac est. Nulla facilisi. Aenean aliquet, dui quis consectetur tempus, nunc elit volutpat nunc, eu euismod arcu diam et arcu. Aenean vulputate metus lectus, id euismod ipsum.4
\\ \\
Vienna, February 12th, 2010  \hfill Julian Schrittwieser

\singlespacing
\tableofcontents
%\onehalfspacing

\chapter{Introduction}
\pagenumbering{arabic} % Roman numerals
\setcounter{page}{1}

\section{Motivation}
mars rover example

\section{Nomenclature}
Throughout this paper, I use \emph{Measure} to refer to a number annotated with its corresponding units. The exact implementation of this annotation depends on the context, some libraries use static type informations, others fields in an object. \emph{Units of Measure} is used to refer to any system that provides such an implementation, including methods to perform arithmetic with \emph{Measure}s. \emph{Units} are the physical units in abbreviated or full form that the Measure is used to represent. Examples include \verb|m/s| for speed or \verb|kg*m/s^2 == N| for force.

\section{Source Code}
The source code for all benchmarks mentioned can be found on GitHub: \url{https://github.com/Mononofu/Units-of-Measure/tree/master/thesis/benchmarks}. This same repository also contains the source for this document (\url{https://github.com/Mononofu/Units-of-Measure/tree/master/thesis}) as well as the implementation of the Units of Measurement system in Scala (\url{https://github.com/Mononofu/Units-of-Measure}). All code is provided under the GNU General Public License, version 3, as is and without warranty of any kind.

\chapter{Related Work}

In principle, there are two distinct ways to implement units of measurement in a programming language: Either statically as part of the type system, or dynamically with objects that know their own unit. Each approach has its own merits and disadvantages, so in the following I examine a few existing implementations in common languages before creating my own solution.

\section{Dynamic Systems}

Some programming languages do not give their users much choice, because they either completely lack static type systems (dynamic languages like Python or Ruby) or their type system is not sophisticated enough for a units of measurement library. In general, those dynamic systems allow access to unit information at runtime but in turn pay a heavy performance price: Each number is stored as a full blown object (see figure \ref{code:naive_java_measure}), not as a simple primitive like a normal number. Thus, all calculations incur additional memory lookups and values take additional space.

\begin{figure}
\begin{minted}{java}
class Measure<T extends java.lang.Number> {
  T      value;
  Unit   unit;
}
\end{minted}
\caption{A naive implementation of a number annotated with its unit in Java. Actual size depends on the number used, but at least 8 bytes are used by the object itself, and further 4/8 bytes (32/64 bit system, resp.) are needed for the reference to the unit.}
\label{code:naive_java_measure}
\end{figure}



\subsection{Java - JScience}

Java has a standardized Units of Measurement API \citep{Units13} with several implementations, both open source and commercial. I will focus on JScience here \citep{Dautelle11}. In addition to support for physical units, JScience also contains libraries for simple mathematical functions, currencies and a linear algebra module. Here, I only focus on parts pertaining to Units of Measurement.

JScience represents Measures using the class \verb/Amount/, but users are free to create their own implementations of \verb/Measurable/. Measures are created with a flexible static helper function (\verb/valueOf/) from a range of different inputs. For details, see the short usage example below in Figure \ref{code:jscience_example}. Essentially, \verb/Amount/ is a wrapper class that tags a number with its corresponding unit - as shown in \ref{code:naive_java_measure} - and provides a set of helper and conversion functions.

\begin{figure}
\begin{minted}{java}
Amount<Mass> m0            = Amount.valueOf(100, POUND);
Amount<ElectricCurrent> m1 = Amount.valueOf("234 mA").to(MICRO(AMPERE));

// m0 =    100 lb
// m1 = 234000 uA

Unit<Mass> u  = m0.getUnit();
Dimension d   = u.getDimension();

\end{minted}
\caption{Example for JScience from the official documentation \citep{Dautelle11}. Usage is straightforward and unit information is still available at runtime.}
\label{code:jscience_example}
\end{figure}

Unit safety is enforced both by the type system at compile time (Measures are generic classes with their units as type parameters) and comparison functions at run time (\verb/equals/ only returns true if both the value and the interval match, etc.).

Unfortunately, the performance hit due to the wrapper classes is quite severe - I experienced a slowdown of around 80 on my machine in a simple benchmark, compared to normal primitives (\verb/int/ and \verb/double/) (see Figure \ref{bench:jscience}).

At the moment, Units of Measurement Systems in Java are only practicable in systems that are not performance critical. Alas, many systems that would benefit most are performance critical - computer game engines, high frequency trading, etc. Therefore, Units of Measurements remain a niche application in Java.

\begin{figure}
\begin{tabular}{lrrr}
method          & int    & double  & JScience \\
\midrule
addition        & 1.0 ns &  1.8 ns    &   149.5 ns \\
multiplication  & 1.1 ns &  1.9 ns    &   162.3 ns
\end{tabular}
\caption{A simple benchmark of JScience compared to primitives data types. JScience is at least 80 times slower on average.}
\label{bench:jscience}
\end{figure}


\subsection{Python - units}

Python is a dynamic language, so the only way to make a system of measurements work is to create a dynamic system that relies on tagged numbers. That's exactly what \emph{units} \citep{Donohue12} does. Usage is very simple, as seen in Figure \ref{code:python_units}. Thanks to operator overloading and the dynamic nature of Python code, units can be introduced into legacy code with only a few changes.

\begin{figure}
\begin{minted}{python}
metre, second = unit('m'), unit('s')                   # define units
print(metre(10) / second(2))                           # 5 m / s
newton = named_unit('N', ['kg', 'm'], ['s', 's'], 1)   # 1 N = 1 kg*m/s^2
units.predefined.define_units()                        # SI units and more
\end{minted}
\caption{\emph{units} makes it very simple to define custom units and to add conversions to existing units, while also including all standard units for immediate use. Calculation happens just like with primitives, thanks to operator overloading.}
\label{code:python_units}
\end{figure}



\begin{figure}
\begin{tabular}{lrrr}
method          & int    & double  & units \\
\midrule
addition        & 151 ns &  164 ns    &    4.136 ns \\
multiplication  & 159 ns &  165 ns    &   22.352 ns
\end{tabular}
\caption{Even though Python is an interpreted language, there is a large difference between primitive numbers, and those provided by the \emph{units} library. The slowdown is especially severe for the multiplication as a new unit has to be generated for the result.}
\label{bench:python_units}
\end{figure}


\section{Static Systems}

Static systems are the complete opposite compared to dynamic systems: They usually do not provide any unit information at runtime, but also do not incur any performance overhead - all units are checked at compile time. To achieve this, they rely on a powerful type system or even compiler plugins, effectively offloading much of the work on the compiler.

\subsection{F\#}

\verb/F#/ is one of only a few languages that have integrated a system of units into the language itself. This allows for easy usage and definition of units since they can rely on special syntax.
\citep{Kennedy08:1}

As usual, \verb/F#/ comes with all common units - both SI and Imperial - predefined \citep{Kennedy08:2}, but it's also easy to interface with code that is not unit aware: Primitive types like \verb/float/ are just a type alias for \verb/float<1>/ (a float with unit dimension).

Thanks to the language support, units also work with generics, including full type (unit) inference. As an example, if we define a simple function \verb/let sqr(x: float<_>) = x*x/, \verb/F#/ will automatically infer that it takes any \verb/float<'u>/ and returns a \verb/float<'u^2>/ \citep{Kennedy08:3}. This works for types of arbitrary complexity, making the programmers life much simpler.

The support for generics also extends to custom types - it's enough to annotate a type parameter with \verb/[<Measure>]/ to make it a unit parameter \citep{Kennedy08:4}. For a example, see Figure \ref{code:fsharp}.


\begin{figure}
\begin{minted}{fsharp}
[<Measure>] type kg                 // define a unit
let gravity = 9.81<m/s^2>           // use units
// generic type
type Vector2< [<Measure>] 'u > = { x: float<'u>, y: float<'u> }
\end{minted}
\caption{Usage examples for the unit system in F\#.}
\label{code:fsharp}
\end{figure}

\subsection{C++ - Boost.Units}

The \verb/Boost.Unit/ library takes advantage of the power of C++ templates in the form of the Boost Metaprogramming Library (MPL).  It relies on the compiler to adequately optimize the code, but then incurs no performance overhead at runtime \citep{Schabel10}.However, since it heavily relies on template metaprogramming it has stringent requirements for the compiler and only works on new, standard compliant compilers.

Variable definitions work as always (\verb/quantity<length> L = 2.0*meters;/), and thanks to operator overloading arithmetic is exactly the same as for primitives.

\verb/Boost.Unit/ is also independent of any single unit systems - functions can be defined generically, for arbitrary unit systems, using just the types of the dimensions. For an example, see Figure \ref{code:boost_units_generic}. Users can then define their own unit systems to interface with already existing library code.

\begin{figure}
\begin{minted}{c++}
template<class System,class Y>
quantity<unit<energy_dimension,System>,Y>
work(quantity<unit<force_dimension,System>,Y> F,
     quantity<unit<length_dimension,System>,Y> dx)
{
    return F*dx;
}
\end{minted}
\caption{The physical definition of work - computed for an arbitrary unit system, from \citep{Schabel10}.}
\label{code:boost_units_generic}
\end{figure}


\subsection{Haskell - Dimensional}
Haskell's \verb/Dimensional/ module \citep{Buckwalter06} completely relies on the type system and is the most elegant of the unit systems presented here. It relies on a single data type (\verb/Dimensional/, see Figure \ref{code:haskell_dimensional}), parameterized with the exponents of different physical dimensions (Figure \ref{code:haskell_dim}). While only integer exponents are possible, rational exponents can be avoided in practice.

\begin{figure}
\begin{minted}{haskell}
newtype (Variant v, Dims d)
      => Dimensional v d a = Dimensional a deriving (Show, Eq, Ord)
\end{minted}
\caption{Dimensional, the unit system's fundamental data type. Since 'a', representing the underlying number, is the only non-phantom type Dimensional is defined as a newtype and does not incur any runtime overhead. From \citep{Buckwalter06}.}
\label{code:haskell_dimensional}
\end{figure}



\begin{figure}
\begin{minted}{haskell}
data (NumType l,    -- Length.
      NumType m,    -- Mass.
      NumType t,    -- Time.
      NumType i,    -- Electric current.
      NumType th,   -- Thermodynamic temperature.
      NumType n,    -- Amount of substance.
      NumType j)    -- Luminous intensity.
   => Dim l m t i th n j

type DLength      = Dim Pos1 Zero Zero Zero Zero Zero Zero
type DForce       = Dim Pos1 Pos1 Neg2 Zero Zero Zero Zero
\end{minted}
\caption{The type Dim, containing the physical base dimensions, and two derived dimesions, from \citep{Buckwalter06}.}
\label{code:haskell_dim}
\end{figure}


As with the other systems, operator overloading allows for easy calculations. \verb/*~/ and \verb|/~| can be used to attach units to numbers and remove them again.


\begin{figure}
\begin{minted}{haskell}
escapeVelocity :: (Floating a) => Mass a -> Length a -> Velocity a
escapeVelocity m r = sqrt (two * g * m / r)
  where
    two = 2 *~ one
    g = 6.6720e-11 *~ (newton * meter ^ pos2 / kilo gram ^ pos2)
\end{minted}
\caption{An example function using the Dimensional module, from \citep{Buckwalter06}.}
\label{code:haskell_dimensional}
\end{figure}

\verb/Dimensional/ is restricted to the physical dimensions mentioned in Figure \ref{code:haskell_dim}. To add more dimensions, the fundamental data type would have to be changed - unfeasible in practice.


\subsection{Scala}

Scala also has a long history of attempts to introduce Units of Measure. One of the oldest of such attempts is a now defunct compiler plugin \citep{Nygard09}. It allowed the programmer to introduce units using multiplication (\verb/400.0*m*m/) and caught common errors when combining units in invalid ways. Sadly, it's only compatible with Scala up to version 2.7.1 - as of 2013, the most recent version is 2.10.0.

Another such example comes from \citep{McBeath08}, who implements a simple unit system based on Church Numerals. However, just as Haskell's \verb/Dimensional/ it is restricted to predefined units and does not allow the user to define her own units.

Finally, there is \verb/ScalaQuantity/ \citep{Hans12}, which is closely modeled on Haskell's \verb/Dimensional/. Again, it does not allow user defined units and is otherwise so similar to \verb/Dimensional/ that I won't cover it in detail.

Currently, there is no Units of Measure system in Scala that allows user defined units or provides for calculations without runtime overhead. This inspired me to create my own library, which I cover in the following.

\chapter{Objectives}

As we have seen, there are a few common features that almost all systems support (custom units, conversions), while some aspects are limited to more advanced libraries (generics, no runtime overhead). I aim to make my Units of Measure system as comprehensive as possible, therefore I present a short summary of all desirable features here. This summary serves as a reference for my implementation, as well as a handy reminder for the reader.


\begin{description}
  \item[No Runtime Overhead] This is one of the most important goals - in many systems that profit from unit checking (games, embedded systems), computing power is at a premium during runtime, while processing at compile time is of little consequence. This also extends to memory consumption - Measures should need no more space than primitive numbers. In fact, they should be indistinguishable at runtime.

  \item[Static Unit Checking] Any mistakes in the units of variables should be detected at compile time - there should not be any surprises at runtime! This is also closely related to the goal above. In fact, it is only possible to check units at compile time, as there is no unit information left at runtime. A simple way to achieve this - also the own I intend to use - is to offload the work of enforcing unit correctness onto the type system, by for example specifying the units as type parameters.

  \item[Custom Units] While it is imperative to provide a thorough selection of units - the SI system at the very least -, it is simply impossible to collect and implement every possible unit. One Application might make heavy use from different types of currency, while another is solely concerned with physical temperatures. Therefore it should be possible for the user to define his own units, as well as to connect them to already existing units so that automatic conversion can take place.

  \item[Unit Conversion] To take full advantage of the provided units, helper methods should automatically convert numbers between different unit systems, as specified by the user. For example, \verb|u(10, "m/s").as("km/h")| automatically changes the units of the Measure and adjusts its value accordingly, to \verb/36/. Naturally, this relationship between units needs to be reflected in their definition, otherwise automatic conversions are impossible.

  \item[Access to a Measure's Units] While it is impossible to access unit information at runtime - see the first goal -, it should still be possible to query the unit of a Measure when its type is known. This is essentially a way to pretty print the underlying type information for a human reader.

  \item[Generics] Another very important aspect is the support of generic classes and methods. A clarifying example: \verb|def square[T](n: T) = n*n| should automatically infer that its return type is T with all units squared. The same principle extends to classes, the most common use case being a simple vector of numbers that defines operations like addition or dot product.

  \item[Implementations for common primitives, as well as all instances of Numeric] Finally, it should be possible to use every type that is an instance of Numeric (that is, every type that can be used as a number) as the underlying type for a Measure. Additionally, special support should be provided for primitives (int, long, float, double), since their performance is the most critical.

\end{description}

\section{Limitations}

Before I describe my implementation in detail, I provide a short list of its limitations. Those are mostly based on limitations of the JVM itself, or on more fundamental issues.

\begin{description}
  \item[Measure is not an instance of Numeric] While Measure requires its underlying type to be an instance of Numeric, it cannot be such an instance itself. The reason is simple: Multiplication will result in a Measure with a different type than the two operands (the units change), thus violating the type bounds of Numeric. The same can even happen for simple addition: Since all operations with a Measure canonicalize (reduce and sort) its type argument (representing its unit), it is possible that an addition of two Measures results in a Measure with a different type.

  \item[Runtime equality ignores units] Since one fundamental goal was to incur no runtime overhead, it is obviously impossible to check unit equality at runtime. In fact, due to the nature of my implementation, Measures based on primitives are indistinguishable from those primitives at runtime - the wrapping Measure class only exists during compilation.

  \item[Generics require Macros] In Java and all languages based on it, generics are implemented in a simple way: The type parameter is only known to the compiler, at runtime the generic class is represented by its raw counterpart. Therefore, all type parameters share the same single method/class - unlike C++, where a separate instance of the generic class is generated for each and every different type parameter. This presents a problem for the Measure macros, since they need to know the Measure's actual type parameter to correctly perform their work (In normal Scala code, it is only known that there is some \verb|Measure[T]|, not very helpful).

  The solution to this problem is to define all generic methods that use Measures as macros, essentially mirroring the way C++ templates are implemented in Scala. As a downside, the Measure macros will be called every time the generic method is used, not just once for each type parameter.

  \item[Arrays lead to degraded Performance] Due to choices made in the current implementation of value classes in Scala, and issues with arrays on the JVM, Measures experience degraded performance when placed in arrays. While the Measure wrapper class is normally only present at compile time for primitives, arrays force the Measure class to stay present - essentially requiring an explicit boxing/unboxing step for each calculation done with the Measure, and also resulting in an increased memory usage. As a workaround, Vector classes with explicit members can be used to reduce the memory overhead.


\end{description}

\chapter{Implementation}

The implementation of my Units of Measure system rests on a few cornerstones that I present here. Without those, this system would not be possible, at least not without sacrificing a few of its goals. Some of those features (like Type Macros) are still very experimental and only available in nightly builds of Scala 2.11. Therefore, my library does not work with the current stable version of Scala (2.10 at the of writing). Luckily, it is very easy to use Scala 2.11 even now - builds are provided over a Sonatype repository and can be used with a simple sbt include (sbt is the de facto standard build system for Scala).

\section{Macros in Scala}

basic principles

\subsection{Context and the Universe}
\subsection{Type Macros}
defining custom units and storing book-keeping information

\section{Value Classes}

achieving 0-overhead

\section{Parser Combinators}

parsing unit expressions from the helper macros and Measure type parameters

\chapter{Usage}

\section{Defining Units}
\section{Arithmetic with Measures}

\chapter{Future Work}

better generics

improvements to value classes


\chapter{Acknowledgements}
scala mailing list

\#scala on freenode

scala user group vienna

Eugene, creator of macros


\singlespacing
\bibliographystyle{dinat}
\bibliography{thesis}{}
\addcontentsline{toc}{chapter}{A \;References}

\chapter*{}
\onehalfspacing
Ich erkläre, dass ich diese Fachbereichsarbeit ausschließlich selbst und ohne Gebrauch unerlaubter Hilfsmittel oder Hilfen verfasst habe.



\end{document}
